\documentclass[a4paper,12pt]{article}

\usepackage{gialosai-style}

\begin{document}

\gialosaiTitlePage{Decision Log}{25 Agosto 2026}{v2.0.0}

\section*{Convenzione}
Ogni decisione \`{e} documentata con: contesto, alternativa scelta, alternative considerate (con confronto tecnico), motivazione della scelta e conseguenze sull'implementazione. Le alternative sono ordinate per vicinanza alla scelta effettiva.

\newpage

% ============================================================
% DC-001
% ============================================================
\begin{decisionbox}
\section*{DC-001 --- Framework per l'interfaccia testuale (TUI)}
\textbf{Decisione.} Il client \`{e} una TUI Go basata su \textbf{Bubble Tea v2} (Model--View--Update) con \textbf{Lipgloss v2} per lo stile.
\end{decisionbox}

\subsection*{Contesto}
Il progetto richiede un'interfaccia utente eseguibile da terminale, senza dipendenze da GUI o browser. Il client deve gestire navigazione multi-pagina, input da tastiera, ridimensionamento terminale e rendering di tabelle con stili colorati.

\gialosaiNote{Il framework Bubble Tea \`{e} utilizzato da strumenti CLI di grande successo come \texttt{lazygit} (25k+ stelle GitHub), \texttt{glow} (rendering markdown), \texttt{soft-serve} (server Git TUI) e \texttt{charm} (la piattaforma Charm). Il pattern MVU (Model-View-Update) \`{e} lo stesso di \texttt{Elm} e \texttt{Redux}, reso popolare da applicazioni web reattive.}

\subsection*{Alternative considerate}

\begin{longtable}{p{0.18\textwidth}p{0.35\textwidth}p{0.35\textwidth}}
\toprule
\tableHeader \textbf{Alternativa} & \textbf{Vantaggi} & \textbf{Svantaggi} \\
\midrule
\endfirsthead
\toprule
\tableHeader \textbf{Alternativa} & \textbf{Vantaggi} & \textbf{Svantaggi} \\
\midrule
\endhead

Bubble Tea (scelta) & Architettura MVU ben definita; ecosistema Charm ricco (Lipgloss, Bubbles, Bubbles/Table); comunit\`{a} attiva; supporto nativo a mouse, AltScreen e window resize; v2 con API stabilizzata; heatingcooling con modello funzionale puro. & Learning curve iniziale per il paradigm MVU; dipendenza dall'ecosistema Charm; v2 ancora recente (possibili breaking change). \\
\midrule

tview & Interfaccia a widget pronti (table, list, form, modal); supporto scroll e mouse; simile ad Android View; documentazione semplice. & Architettura ad oggetti con mutabilit\`{a} diffusa; meno flessibile per layout personalizzati; stile visivo pi\`{u} datato; comunit\`{a} pi\`{u} piccola; nessun sistema di styling separato. \\
\midrule

termui & Leggero, senza dipendenze esterne; API minimale; buono per dashboard semplici. & Nessun sistema di layout dichiarativo; gestione manuale dello stato; nessun supporto a widget complessi (tabelle, form); manutenzione ridotta. \\
\midrule

\texttt{os/exec} + input raw & Nessuna dipendenza; controllo totale su byte-level del terminale. & Implementazione manuale di ANSI escape codes; nessuna astrazione di layout; testabilit\`{a} quasi nulla; codice non manutenibile. \\
\bottomrule
\end{longtable}

\subsection*{Motivazione}
Bubble Tea \`{e} la scelta naturale per TUI complesse in Go: il paradigm MVU separa stato, logica e rendering; Lipgloss fornisce un sistema di styling dichiarativo; Bubbles offre componenti pronti (table, viewport, text input). L'ecosistema Charm ha una comunit\`{a} attiva e documentazione completa. tview avrebbe limitato la personalizzazione visiva; termui e os/exec avrebbero richiesto troppo lavoro manuale.

\gialosaiNote{Esempi concreti di Bubble Tea in produzione: \texttt{lazygit} (interfaccia per git), \texttt{k9s} (amministrazione Kubernetes), \texttt{btop} (monitor risorse), \texttt{glow} (rendering markdown), \texttt{soft-serve} (server Git TUI), \texttt{mods} (generazione testo con IA). Questi dimostrano la maturit\`{a} e l'affidabilit\`{a} del framework.}

\subsection*{Conseguenze}
Tutte le pagine seguono il pattern \texttt{Init/Update/View}. Lo stile \`{e} centralizzato in \texttt{client/ui/styles/}. La navigazione tra pagine \`{e} gestita dal modello principale \texttt{mainLayout} con enum \texttt{iota}. Il tabnav usa il rendering reale per calcolare le hitbox dei click.

% ============================================================
% DC-002
% ============================================================
\begin{decisionbox}
\section*{DC-002 --- Router HTTP per il server}
\textbf{Decisione.} Il server usa \textbf{chi} (v5.3.0) come router HTTP.
\end{decisionbox}

\subsection*{Contesto}
Il backend Go deve esporre un insieme di rotte REST organizzate per area di dominio, con separazione tra rotta pubblica (login) e rotte protette (middleware auth). Il router deve supportare middleware a catena, parametri URL, gruppi di rotte e method matching.

\gialosaiNote{chi \`{e} utilizzato da progetti come \texttt{go-chi/chi} (naturalmente), \texttt{PresslyLabs/goose} e molti altri. Il design di chi \`{e} ispirato a \texttt{net/http} standard, garantendo compatibilit\`{a} con l'intero ecosistema Go.}

\subsection*{Alternative considerate}

\begin{longtable}{p{0.16\textwidth}p{0.35\textwidth}p{0.35\textwidth}}
\toprule
\tableHeader \textbf{Alternativa} & \textbf{Vantaggi} & \textbf{Svantaggi} \\
\midrule
\endfirsthead
\toprule
\tableHeader \textbf{Alternativa} & \textbf{Vantaggi} & \textbf{Svantaggi} \\
\midrule
\endhead

chi (scelta) & Interfaccia compatibile con \texttt{net/http} standard; middleware chain; route groups; URL params; stile composable; nessuna dipendenza esterna pesante; buona documentazione. & Nessun built-in validation; routing tree radix-3 leggermente pi\`{u} lento di gin in benchmark estremi (differenza trascurabile). \\
\midrule

gin-gonic & Performance estremamente alte (radix tree); binding JSON integrato; middleware built-in; validazione integrata; community grande. & Interfaccia \texttt{*gin.Context} non \`{e} \texttt{http.Request} nativo; accoppiamento stretto con il framework; package routing pi\`{u} rigido; meno trasparente nei middleware. \\
\midrule

echo & Performance alte; middleware chain; routing con group; validazione; OpenAPI automatico. & \texttt{*echo.Context} wrapper proprietario; meno trasparente di chi per il passaggio di contesto; community leggermente pi\`{u} piccola di gin/chi. \\
\midrule

\texttt{net/http} puro & Nessuna dipendenza; controllo totale; stdlib. & Routing manuale con \texttt{ServeMux} (limitato prima di Go 1.22); nessun supporto a parametri URL nativo; codice boilerplate elevato; nessun middleware chain organizzato. \\
\bottomrule
\end{longtable}

\subsection*{Motivazione}
chi \`{e} un wrapper leggero di \texttt{net/http}: le rotte sono handler standard, i middleware sono \texttt{func(http.Handler) http.Handler}, e il contesto \`{e} quello nativo di Go. Questo garantisce trasparenza, compatibilit\`{a} con l'ecosistema e facilit\`{a} di testing con \texttt{httptest}. gin sarebbe stato pi\`{u} performante ma pi\`{u} invasivo; echo avrebbe aggiunto un wrapper inutile; \texttt{net/http} puro avrebbe richiesto troppo boilerplate per il routing con parametri.

\gialosaiNote{Il pattern di chi \`{e} lo stesso utilizzato da \texttt{docker} (che usa \texttt{mux} di stdlib) e \texttt{kubectl} (che usa \texttt{cobra} per CLI e \texttt{net/http} per API). La scelta di chi rappresenta un compromesso ottimale tra semplicit\`{a} e funzionalit\`{a}.}

\subsection*{Conseguenze}
Il router \`{e} composto in \texttt{router.NewRouter} dove vengono creati tutti i repository, service e controller. Le rotte protette usano \texttt{r.Group()} con \texttt{middleware.NewAuthMiddleware(sessionRepo)}. I parametri URL (es. \texttt{\{username\}}, \texttt{\{order\_id\}}) vengono estratti con \texttt{chi.URLParam}.

% ============================================================
% DC-003
% ============================================================
\begin{decisionbox}
\section*{DC-003 --- Database relazionale}
\textbf{Decisione.} Il sistema usa \textbf{MySQL 8.0} eseguito tramite Docker Compose, con driver \texttt{go-sql-driver/mysql}.
\end{decisionbox}

\subsection*{Contesto}
Il sistema richiede un database relazionale condiviso multiutente, vincoli di integrit\`{a} referenziale, transazioni e supporto a operazioni multi-entit\`{a} (ordini, righe, giacenze, DDT, fatture).

\gialosaiNote{MySQL 8 \`{e} lo stesso database utilizzato da \texttt{docker} (per Docker Hub), \texttt{kubernetes} (per gli store) e \texttt{github} (per GitHub.com). Il driver \texttt{go-sql-driver/mysql} \`{e} il pi\`{u} utilizzato nell'ecosistema Go per MySQL.}

\subsection*{Alternative considerate}

\begin{longtable}{p{0.16\textwidth}p{0.35\textwidth}p{0.35\textwidth}}
\toprule
\tableHeader \textbf{Alternativa} & \textbf{Vantaggi} & \textbf{Svantaggi} \\
\midrule
\endfirsthead
\toprule
\tableHeader \textbf{Alternativa} & \textbf{Vantaggi} & \textbf{Svantaggi} \\
\midrule
\endhead

MySQL 8 (scelta) & InnoDB con ACID; chiavi esterne; \texttt{utf8mb4}; Docker image ufficiale; healthcheck; ecosistema maturo; MySQL 8 supporta CTE e window functions. & Overhead rispetto a SQLite per sviluppo singolo; configurazione iniziale pi\`{u} complessa;occupa pi\`{u} risorse. \\
\midrule

SQLite & Zero configurazione; file singolo; ideale per prototipi; embeddabile. & Nessun supporto concorrente multi-writer; nessun controllo di accesso multi-utente; funzionalit\`{a} SQL ridotte; nessun supporto nativo a stored procedures; InnoDB perde vincoli di integrit\`{a}. \\
\midrule

PostgreSQL & Funzionalit\`{a} SQL pi\`{u} ricche (JSON, array, CTE avanzate); MVCC nativo; estensioni (PostGIS, pg\_trgm). & Immagine Docker pi\`{u} pesante; configurazione iniziale pi\`{u} complessa di MySQL; per un progetto universitario le funzionalit\`{a} extra non sono necessarie. \\
\midrule

MongoDB & Schema flessibile; scalabilit\`{a} orizzontale; JSON nativo. & Nessuna relazione nativa; nessuna transazione multi-documento garanzia (prima di v4.0); inadatto a dominio relazionale come ordini/giacenze; perdita di vincoli di integrit\`{a}. \\
\bottomrule
\end{longtable}

\subsection*{Motivazione}
MySQL 8 con InnoDB \`{e} il compromesso ottimale: fornisce ACID, chiavi esterne, transazioni e \texttt{utf8mb4} con un'immagine Docker leggera e ben supportata. SQLite non supporta concorrenza multi-writer; PostgreSQL avrebbe aggiunto complessit\`{a} non necessaria; MongoDB \`{e} inadatto a un dominio relazionale con vincoli di integrit\`{a} forti (ogni prodotto ha una giacenza, le righe dipendono dagli ordini).

\subsection*{Conseguenze}
Lo schema \`{e} definito in \texttt{docker/mysql/init.sql} con InnoDB e \texttt{utf8mb4}. Docker Compose avvia MySQL sulla porta host 3307 con healthcheck. I repository usano \texttt{database/sql} con query parametrizzate. Le transazioni sono gestite da \texttt{TransactionHelper} che propagates nel contesto.

% ============================================================
% DC-004
% ============================================================
\begin{decisionbox}
\section*{DC-004 --- Architettura a layer con Repository pattern}
\textbf{Decisione.} Il server \`{e} organizzato in layer: \textbf{router $\to$ controller $\to$ service $\to$ repository $\to$ domain}, con \textbf{Repository pattern} per l'accesso ai dati e dipendenze iniettate via interfacce.
\end{decisionbox}

\subsection*{Contesto}
Le responsabilit\`{a} di trasporto HTTP, logica applicativa e persistenza SQL devono restare separate e indipendenti. I service devono essere testabili con mock senza accedere al database.

\gialosaiNote{Il Repository pattern \`{e} lo stesso utilizzato da \texttt{database/sql} in Go, da \texttt{gorm} e da \texttt{sqlx}. L'architettura a layer \`{e} lo standard per backend Go, utilizzato da \texttt{docker} (che usa un pattern simile con client/API/server) e \texttt{kubernetes} (che usa un pattern simile con informer/cache/store).}

\subsection*{Alternative considerate}

\begin{longtable}{p{0.16\textwidth}p{0.35\textwidth}p{0.35\textwidth}}
\toprule
\tableHeader \textbf{Alternativa} & \textbf{Vantaggi} & \textbf{Svantaggi} \\
\midrule
\endfirsthead
\toprule
\tableHeader \textbf{Alternativa} & \textbf{Vantaggi} & \textbf{Svantaggi} \\
\midrule
\endhead

Layer + Repository (scelta) & Separazione chiara delle responsabilit\`{a}; testabilit\`{a} con mock; coerenza con Dependency Inversion; ogni layer \`{e} sostituibile; modello standard per backend Go. & Pi\`{u} boilerplate (interfacce, iniezione); crescita del numero di file; rischio di ``service anemici'' se la logica si disperde. \\
\midrule

MVC classico & Pattern noto; separazione model/view/controller. & La view \`{e} nel client separato; il controller HTTP non \`{e} un MVC; il pattern non si adatta alla separazione client/server. \\
\midrule

DAO generico & Accesso ai dati pi\`{u} flessibile; meno file. & Accoppia i service al dialetto SQL; nessuna separazione tra logica e persistenza; difficile testare con mock. \\
\midrule

Moduli monolitici (senza layer) & Zero boilerplate; codice compatto. & Test impossibili senza DB; logica mescolata a SQL; impossibile sostituire componenti; debito tecnico crescente. \\
\bottomrule
\end{longtable}

\subsection*{Motivazione}
Il Repository pattern isola il dialetto SQL nei repository, che espongono interfacce ristrette. I service dipendono da queste interfacce (Dependency Inversion), non dalle implementazioni concrete. Questo consente di sostituire i repository con mock nei test, come effettivamente fatto con \texttt{go-sqlmock}. MVC non si adatta perch\'{e} la view \`{e} nel processo client separato. DAO non fornirebbe la separazione tra logica e persistenza.

\gialosaiNote{Il pattern \`{e} identico a quello utilizzato da \texttt{docker} (che separa client/API/daemon) e \texttt{kubernetes} (che separa API server/controller/scheduler/node). La separazione garantisce testabilit\`{a} e manutenibilit\`{a}.}

\subsection*{Conseguenze}
Ogni area di dominio (\texttt{auth}, \texttt{user}, \texttt{product}, \texttt{category}, \texttt{order}) ha i propri: controller (adattatore HTTP), service (logica di business), repository (accesso SQL) e DTO (rappresentazione di trasporto). Il router compone le dipendenze in \texttt{NewRouter}. I test dei service usano mock delle interfacce di repository.

% ============================================================
% DC-005
% ============================================================
\begin{decisionbox}
\section*{DC-005 --- Autenticazione JWT con sessione persistita}
\textbf{Decisione.} Il server usa \textbf{JWT} (HS256, \texttt{golang-jwt/jwt/v5}) con una \textbf{whitelist di sessioni} nel database (\texttt{Sessioni}). Il token viene memorizzato come hash SHA-256, mai in chiaro.
\end{decisionbox}

\subsection*{Contesto}
Il sistema deve autenticare le richieste HTTP, supportare logout, gestire la scadenza delle sessioni e proteggere da token rubati.

\gialosaiNote{L'autenticazione JWT con whitelist \`{e} lo stesso approccio utilizzato da \texttt{docker} (con token di accesso) e \texttt{kubectl} (con certificati e token). Il pattern di hash SHA-256 del token \`{e} una best practice per evitare che token compromessi siano riutilizzabili.}

\subsection*{Alternative considerate}

\begin{longtable}{p{0.16\textwidth}p{0.35\textwidth}p{0.35\textwidth}}
\toprule
\tableHeader \textbf{Alternativa} & \textbf{Vantaggi} & \textbf{Svantaggi} \\
\midrule
\endfirsthead
\toprule
\tableHeader \textbf{Alternativa} & \textbf{Vantaggi} & \textbf{Svantaggi} \\
\midrule
\endhead

JWT + sessione DB (scelta) & Logout effettivo (revoca lato server); controllo della scadenza lato DB; token memorizzato come hash; sicurezza contro token rubati. & Overhead di una query per ogni richiesta protetta; persistenza della sessione nel DB; complessit\`{a} di gestione della tabella Sessioni. \\
\midrule

JWT stateless (senza DB) & Nessuna query per richiesta; performance migliori; scalabilit\`{a}. & Impossibile fare logout effettivo; token valido fino alla scadenza anche se rubato; nessun controllo server-side sulla sessione; revoca impossibile. \\
\midrule

Sessioni classiche (cookie + DB) & Standard web; revoca semplice; non serve gestire JWT. & Dipendenza dai cookie; non \`{e} lo standard per API REST; problemi di CSRF; il client TUI deve gestire cookie manualmente. \\
\midrule

OAuth2 / OpenID Connect & Standard industriale; federazione; SSO. & Overkill per un sistema locale con 3 ruoli; infrastruttura complessa; dipendenza da provider esterno. \\
\bottomrule
\end{longtable}

\subsection*{Motivazione}
JWT stateless non consente il logout effettivo: se il token viene rubato, resta valido fino alla scadenza. Con la whitelist nel database, il middleware verifica sia la firma sia la presenza della sessione non scaduta nella tabella \texttt{Sessioni}. L'hash SHA-256 del token viene memorizzato, cos\`{i} che anche una compromissione del database non riveli i token. Sessioni classiche con cookie non sono adatte a un client TUI.

\gialosaiNote{Il pattern \`{e} identico a quello utilizzato da \texttt{docker} (con token di accesso) e \texttt{kubectl} (con certificati e token). La whitelist nel database \`{e} la best practice per sistemi che richiedono logout effettivo e controllo server-side della sessione.}

\subsection*{Conseguenze}
Il login genera un JWT con claim \texttt{user\_id}, \texttt{username}, \texttt{role} e scadenza 24h. Il token viene hashato con SHA-256 e salvato nella tabella \texttt{Sessioni}. Il middleware per ogni richiesta protetta: (1) verifica la firma JWT, (2) hash il token, (3) consulta il DB per validit\`{a} della sessione, (4) inietta \texttt{userID} nel contesto. Il logout elimina la riga dalla tabella.

% ============================================================
% DC-006
% ============================================================
\begin{decisionbox}
\section*{DC-006 --- Password hashing con bcrypt}
\textbf{Decisione.} Le password sono hashate con \textbf{bcrypt} (\texttt{golang.org/x/crypto/bcrypt}, costo di default 10).
\end{decisionbox}

\subsection*{Contesto}
Le password devono essere memorizzate in modo sicuro, senza testo in chiaro, resistendo ad attacchi a dizionario e rainbow table.

\gialosaiNote{bcrypt \`{e} lo standard utilizzato da \texttt{docker} (per le password degli utenti), \texttt{kubernetes} (per le secret) e \texttt{github} (per le password degli utenti). Il costo di default (10) \`{e} lo stesso utilizzato dalla maggior parte dei sistemi di autenticazione.}

\subsection*{Alternative considerate}

\begin{longtable}{p{0.16\textwidth}p{0.35\textwidth}p{0.35\textwidth}}
\toprule
\tableHeader \textbf{Alternativa} & \textbf{Vantaggi} & \textbf{Svantaggi} \\
\midrule
\endfirsthead
\toprule
\tableHeader \textbf{Alternativa} & \textbf{Vantaggi} & \textbf{Svantaggi} \\
\midrule
\endhead

bcrypt (scelta) & Lento per design (costo configurabile); resistente a rainbow table; salt integrato; standard industriale; libreria Go ufficiale (\texttt{x/crypto}). & Costo fisso di CPU (trascurabile per login su rete locale); hash lungo (60 byte). \\
\midrule

scrypt & Parametri configurabili (N, r, p); resistente a GPU/ASIC; usato in cryptocurrencies. & Parametri da calibrare; meno maturo di bcrypt;implementazioni Go meno note. \\
\midrule

argon2id & State-of-the-art ( Winner Password Hashing Competition); resistente a side-channel; parametri configurabili. & Pi\`{u} recente; meno supportato in librerie esterne; parametri da calibrare; ecosistema Go meno maturo. \\
\midrule

SHA-256 (nonsicuro) & Velocissimo; libreria stdlib. & Vulnerabile a rainbow table; nessun salt; non \`{e} un password hasher; reversibile con dictionary attack. \\
\bottomrule
\end{longtable}

\subsection*{Motivazione}
bcrypt \`{e} lo standard consolidato per il hashing delle password. Il costo di default (10) \`{e} sufficiente per un sistema locale. SHA-256 non \`{e} sicuro per password. scrypt e argon2id sono pi\`{u} moderni ma per un progetto universitario la differenza \`{e} trascurabile; bcrypt \`{e} pi\`{u} collaudato e documentato.

\subsection*{Conseguenze}
\texttt{HashPassword} genera l'hash con \texttt{bcrypt.GenerateFromPassword} e costo 10. \texttt{CheckPassword} verifica con \texttt{bcrypt.CompareHashAndPassword}. L'hash \`{e} memorizzato nella colonna \texttt{password\_hash VARCHAR(255)} della tabella \texttt{Utenti}.

% ============================================================
% DC-007
% ============================================================
\begin{decisionbox}
\section*{DC-007 --- TransactionHelper con contesto}
\textbf{Decisione.} \texttt{domain.TransactionHelper} inserisce una transazione SQL nel \texttt{context.Context} della richiesta; i repository la estraggono e la usano quando disponibile, altrimenti aprono una connessione diretta tramite l'interfaccia \texttt{QueryExecutor}.
\end{decisionbox}

\subsection*{Contesto}
Operazioni come registrazione ordini, aggiornamento giacenze e generazione DDT richiedono atomicit\`{a} su pi\`{u} tabelle. La transazione deve essere trasparente ai service e compatibile con il pattern dei repository.

\gialosaiNote{Il pattern di transazione in contesto \`{e} lo stesso utilizzato da \texttt{database/sql} di Go e da ORM come \texttt{gorm}. La propagazione via \texttt{context.Context} \`{e} la best practice raccomandata dal team di Go per la gestione delle transazioni.}

\subsection*{Alternative considerate}

\begin{longtable}{p{0.16\textwidth}p{0.35\textwidth}p{0.35\textwidth}}
\toprule
\tableHeader \textbf{Alternativa} & \textbf{Vantaggi} & \textbf{Svantaggi} \\
\midrule
\endfirsthead
\toprule
\tableHeader \textbf{Alternativa} & \textbf{Vantaggi} & \textbf{Svantaggi} \\
\midrule
\endhead

Transaction-in-context (scelta) & Trasparente ai repository; rollback automatico; supporta nesting; i service non gestiscono \texttt{BEGIN/COMMIT} manualmente; testabile. & Uso di \texttt{context.WithValue} (type-safe only at boundary); debugging pi\`{u} complesso; invisibilit\`{a} implicita della transazione. \\
\midrule

Manuale (\texttt{tx, err := db.Begin()}) & Esplicito; facile da capire; debugging semplice. & Boilerplate in ogni service; rischio di dimenticare rollback; codice duplicato; difficile nesting. \\
\midrule

ORM (GORM) & API ad alto livello; migration; callback pre/post. & Abstraction leak; SQL meno controllabile; dipendenza pesante; overhead per un progetto con schema noto. \\
\bottomrule
\end{longtable}

\subsection*{Motivazione}
Il pattern transaction-in-context \`{e} il pi\`{u} idoneo per un backend Go con repository pattern: il service chiama \texttt{Tx.WithinTransaction(ctx, fn)} e tutti i repository all'interno della closure partecipano alla stessa transazione, senza sapere di essere in una transazione. L'interfaccia \texttt{QueryExecutor} (che unisce \texttt{ExecContext}, \texttt{QueryContext}, \texttt{QueryRowContext}) \`{e} implementata sia da \texttt{*sql.Tx} che da \texttt{*sql.DB}, rendendo la transizione trasparente.

\subsection*{Conseguenze}
\texttt{TransactionHelper.WithinTransaction} verifica se esiste gi\`{a} una transazione nel contesto (supporto nesting), altrimenti ne crea una con \texttt{DB.BeginTx}. In caso di errore esegue rollback; in caso di successo commit. I repository chiamano \texttt{GetExecutor(ctx, r.DB)} per ottenere l'executor corretto. Il meccanismo \`{e} usato in \texttt{UserService\_Login} (atomicit\`{a} di read-modify-write per tentativi falliti) e nella creazione ordini (aggiornamento atomico di ordini e giacenze).

% ============================================================
% DC-008
% ============================================================
\begin{decisionbox}
\section*{DC-008 --- Nomi delle pagine con iota}
\textbf{Decisione.} I nomi delle pagine della TUI sono definiti come costanti \texttt{iota} in un enum Go, anzich\'{e} come stringhe o interi manuali.
\end{decisionbox}

\subsection*{Contesto}
Il router del client deve commutare tra decine di viste (login, magazzino, contabilit\`{a}, admin e sottopagine) in base al valore dello stato. La navigazione avviene tramite \texttt{switch} nel metodo \texttt{Update} e \texttt{View} del modello principale.

\gialosaiNote{Il pattern di enum con \`{e} lo stesso utilizzato da \texttt{lazygit}, \texttt{k9s} e \texttt{btop} per la navigazione. Il sistema di navigazione a enum garantisce che il compilatore Go possa rilevare switch incompleti, aiutando a identificare pagine non gestite.}

\subsection*{Alternative considerate}

\begin{longtable}{p{0.16\textwidth}p{0.35\textwidth}p{0.35\textwidth}}
\toprule
\tableHeader \textbf{Alternativa} & \textbf{Vantaggi} & \textbf{Svantaggi} \\
\midrule
\endfirsthead
\toprule
\tableHeader \textbf{Alternativa} & \textbf{Vantaggi} & \textbf{Svantaggi} \\
\midrule
\endhead

iota enum (scelta) & Unicit\`{a} garantita a compile-time; nessun errore di battitura; rende il codice leggibile (\texttt{LOGIN}, \texttt{MAGMENU}); il compilatore avvisa per switch incompleti. & Indici impliciti (meno trasparente in debug); non autoserializzabile in JSON. \\
\midrule

Stringhe costanti & Leggibilit\`{a} diretta; debug semplice. & Possibili duplicati; errori di battitura non rilevati; pi\`{u} lento nel confronto. \\
\midrule

Interi manuali & Performance massime. & Errori di assegnazione non rilevati; codice poco leggibile; rischio di indici duplicati. \\
\midrule

Registry dinamico (map) & Flessibilit\`{a} a runtime; registrazioni dinamiche. & Overhead; nessun controllo a compile-time; complessit\`{a} non necessaria per un enum fisso. \\
\bottomrule
\end{longtable}

\subsection*{Motivazione}
L'uso di \texttt{iota} garantisce unicit\`{a} a compile-time e rende il codice leggibile. Il compilatore Go può rilevare switch incompleti (escluso con \texttt{default}), aiutando a identificare pagine non gestite. Stringhe costanti sarebbero state ugualmente valide ma senza il supporto del compilatore per switch.

\subsection*{Conseguenze}
L'enum \texttt{paginaAttiva} definisce \texttt{LOGIN}, \texttt{MAGMENU}, \texttt{CONTMENU}, \texttt{ADMINMENU}. L'enum \texttt{userRole} definisce \texttt{MAGAZZINIERE}, \texttt{CONTABILE}, \texttt{ADMIN}. Il metodo \texttt{parseUserRole} in \texttt{app.go} converte la stringa del ruolo (ricevuta dal JWT) nel corrispondente valore enum.

% ============================================================
% DC-009
% ============================================================
\begin{decisionbox}
\section*{DC-009 --- Containerizzazione con Docker Compose}
\textbf{Decisione.} Il sistema \`{e} containerizzato con \textbf{Docker Compose} (due servizi: \texttt{db} per MySQL e \texttt{app} per il backend). Il Dockerfile del server usa un \textbf{multi-stage build} con builder Alpine e runtime distroless nonroot.
\end{decisionbox}

\subsection*{Contesto}
Il server richiede MySQL e deve essere avviabile in modo riproducibile. Il database deve essere persistente e inizializzato con script SQL alla prima creazione.

\gialosaiNote{L'approccio Docker segue le best practice di \texttt{docker} e \texttt{docker-compose}: un file YAML che definisce i servizi, le dipendenze e i volumi. Il pattern multi-stage build \`{e} lo stesso utilizzato da \texttt{golang} official image e da progetti come \texttt{traefik}.}

\subsection*{Alternative considerate}

\begin{longtable}{p{0.16\textwidth}p{0.35\textwidth}p{0.35\textwidth}}
\toprule
\tableHeader \textbf{Alternativa} & \textbf{Vantaggi} & \textbf{Svantaggi} \\
\midrule
\endfirsthead
\toprule
\tableHeader \textbf{Alternativa} & \textbf{Vantaggi} & \textbf{Svantaggi} \\
\midrule
\endhead

Docker Compose (scelta) & Avvio con un comando; dipendenze dichiarate; healthcheck; volumi persistenti; script di init automatici; riproducibile. & Overhead di container; debugging pi\`{u} complesso; consumo risorse. \\
\midrule

Avvio manuale (host) & Zero overhead; debugging diretto. & Installazione MySQL manuale; inizializzazione schema manuale; non riproducibile su altre macchine. \\
\midrule

Kubernetes & Orchestrazione scalabile; self-healing; rollbacks. & Complessit\`{a} enorme per un progetto universitario; necessita di cluster; overkill. \\
\midrule

Podman + Podman Compose & Alternativa rootless a Docker; compatibilit\`{a} Docker. & Community pi\`{u} piccola; meno tooling; meno supporto da parte delle CI/CD. \\
\bottomrule
\end{longtable}

\subsection*{Motivazione}
Docker Compose \`{e} lo standard per lo sviluppo locale multi-container. Un singolo comando (\texttt{docker compose up --build}) avvia MySQL con healthcheck e il backend. Gli script di init (\texttt{init.sql}, \texttt{prova.sql}) vengono eseguiti automaticamente alla prima creazione del volume. Il multi-stage build riduce l'immagine finale a un binario statico su distroless nonroot, minimizzando la superficie di attacco.

\subsection*{Conseguenze}
Il \texttt{docker-compose.yml} definisce i servizi \texttt{db} (MySQL 8, porta 3307) e \texttt{app} (Go backend, porta 8080). Il volume \texttt{mysql\_data} persiste i dati. Il Dockerfile usa \texttt{golang:1.25-alpine} per il build e \texttt{gcr.io/distroless/static-debian12:nonroot} per il runtime. Il servizio \texttt{app} dipende da \texttt{db} con \texttt{condition: service\_healthy}.

% ============================================================
% DC-010
% ============================================================
\begin{decisionbox}
\section*{DC-010 --- Scanner come emulazione tastiera}
\textbf{Decisione.} Il lettore codici a barre EAN-13 \`{e} integrato come emulazione di tastiera: il terminale riceve la stringa come input sequenziale di caratteri, senza API dedicate o driver specifici.
\end{decisionbox}

\subsection*{Contesto}
La TUI non pu\`{o} accedere a dispositivi USB direttamente. Il lettore fisico opera come un normale terminale, inviando i caratteri del codice come keystrokes.

\gialosaiNote{L'approccio di emulazione tastiera \`{e} lo stesso utilizzato da \texttt{htop}, \texttt{btop} e \texttt{lazygit} per l'input da tastiera. Il pattern di input come keystroke \`{e} lo standard per applicazioni TUI che devono interoperare con dispositivi fisici.}

\subsection*{Alternative considerate}

\begin{longtable}{p{0.16\textwidth}p{0.35\textwidth}p{0.35\textwidth}}
\toprule
\tableHeader \textbf{Alternativa} & \textbf{Vantaggi} & \textbf{Svantaggi} \\
\midrule
\endfirsthead
\toprule
\tableHeader \textbf{Alternativa} & \textbf{Vantaggi} & \textbf{Svantaggi} \\
\midrule
\endhead

Emulazione tastiera (scelta) & Nessun driver; trasparente alla TUI; funziona con qualsiasi lettore compatibile; zero configurazione. & Dipende dalla velocit\`{a} del lettore; rischio di input parziali se il terminale \`{e} lento; nessun feedback acustico/visivo integrato. \\
\midrule

API USB (libusb) & Controllo diretto del dispositivo; indipendente dal terminale. & Necessita di permessi USB; dipende dal modello di lettore; complessit\`{a} elevata; non portable. \\
\midrule

Protocollo seriale & Comunicazione diretta; feedack immediato. & Configurazione porta seriale; dipendenza dal modello; driver specifici. \\
\midrule

Input manuale da tastiera & Zero dipendenze. & Errore umano elevato; lento; non automatizzabile. \\
\bottomrule
\end{longtable}

\subsection*{Motivazione}
La maggior parte dei lettori codici a barre opera come ``keyboard wedge'': inviano il codice come se fosse stato digitato da tastiera. Questo approccio \`{e} universale, non richiede driver e funziona immediatamente con la TUI. Il codice EAN-13 viene validato dal service e usato per identificare il prodotto.

\subsection*{Conseguenze}
La pagina scanner in \texttt{client/ui/pages/magazzino/scanner/} riceve l'input come normale keystroke. Il codice EAN-13 viene validato con \texttt{domain.IsValidEAN13} prima di essere usato per la ricerca del prodotto. Non \`{e} necessario un driver o un protocollo specifico.

% ============================================================
% DC-011
% ============================================================
\begin{decisionbox}
\section*{DC-011 --- Blocco account dopo 3 tentativi}
\textbf{Decisione.} Il servizio utenti blocca automaticamente l'account dopo \textbf{3 tentativi di login falliti} consecutivi. L'account con username \texttt{admin} \`{e} esente dal blocco. Il ripristino avviene solo tramite un altro admin.
\end{decisionbox}

\subsection*{Contesto}
Il sistema deve proteggersi da attacchi brute-force su una rete locale dove i tentativi di accesso non autorizzati sono un rischio concreto.

\gialosaiNote{Il blocco dopo tentativi falliti \`{e} lo stesso approccio utilizzato da \texttt{ssh} (con \texttt{fail2ban}), \texttt{docker} (con rate limiting) e \texttt{github} (con blocco account). Il pattern \`{e} standard per sistemi che richiedono sicurezza su rete locale.}

\subsection*{Alternative considerate}

\begin{longtable}{p{0.16\textwidth}p{0.35\textwidth}p{0.35\textwidth}}
\toprule
\tableHeader \textbf{Alternativa} & \textbf{Vantaggi} & \textbf{Svantaggi} \\
\midrule
\endfirsthead
\toprule
\tableHeader \textbf{Alternativa} & \textbf{Vantaggi} & \textbf{Svantaggi} \\
\midrule
\endhead

Blocco dopo 3 tentativi (scelta) & Protezione concreta contre-force; semplice da implementare; comportamento prevedibile. & Blocco accidentale legittimo; dipendenza dall'admin per il ripristino; timing attack possibile (ma non critico su rete locale). \\
\midrule

Rate limiting (es. 5 req/min) & Pi\`{u} flessibile; non blocca l'account; grace period. & Pi\`{u} complesso da implementare; non protegge da attacchi lenti; necessita di storage dei tentativi. \\
\midrule

CAPTCHA & Protezione contro bot; non blocca l'account. & Non applicabile in TUI; complessit\`{a} di integrazione;esperienza utente ridotta. \\
\midrule

Nessuna protezione & Zero overhead. & Vulnerabile a brute-force; inaccettabile per un sistema con dati sensibili. \\
\bottomrule
\end{longtable}

\subsection*{Motivazione}
Per un sistema locale con 3 ruoli, il blocco dopo 3 tentativi \`{e} la soluzione pi\`{u} pratica: semplice, efficace e documentata. L'esenzione dell'admin garantisce sempre l'accesso amministrativo. Il rate limiting sarebbe stato pi\`{u} flessibile ma pi\`{u} complesso; CAPTCHA non \`{e} applicabile in TUI.

\subsection*{Conseguenze}
La costante \texttt{MaxTentativiLogin = 3} \`{e} definita in \texttt{domain/constants.go}. Il metodo \texttt{UserService\_Login} gestisce il ciclo: (1) fetch utente, (2) verifica admin, (3) verifica blocco, (4) verifica password, (5) incrementa tentativi o resetta. Tutto avviene in un'unica transazione. Dopo 3 fallimenti, \texttt{Bloccato} viene impostato a \texttt{true} e l'utente non pu\`{o} pi\`{u} autenticarsi.

% ============================================================
% DC-012
% ============================================================
\begin{decisionbox}
\section*{DC-012 --- UML come sorgente unico delle immagini}
\textbf{Decisione.} I diagrammi UML sono mantenuti come sorgenti PlantUML in \texttt{documentazione/uml/}; le immagini PNG in \texttt{documentazione/img/} sono esportazioni generate da questi sorgenti.
\end{decisionbox}

\subsection*{Contesto}
Le figure dei documenti LaTeX devono restare allineate ai sorgenti e al codice ed essere rigenerabili in modo riproducibile.

\gialosaiNote{PlantUML \`{e} lo standard per diagrammi UML in formato testuale, utilizzato da \texttt{docker}, \texttt{kubernetes} e \texttt{github} per la documentazione. Il pattern di sorgente testuale + export \`{e} lo stesso utilizzato da \texttt{mermaid} e \texttt{d2}.}

\subsection*{Alternative considerate}

\begin{longtable}{p{0.16\textwidth}p{0.35\textwidth}p{0.35\textwidth}}
\toprule
\tableHeader \textbf{Alternativa} & \textbf{Vantaggi} & \textbf{Svantaggi} \\
\midrule
\endfirsthead
\toprule
\tableHeader \textbf{Alternativa} & \textbf{Vantaggi} & \textbf{Svantaggi} \\
\midrule
\endhead

PlantUML (scelta) & Notazione testuale; versionabile con Git; rigenerabile; extention VS Code; standard per diagrammi UML. & Rendering lento per diagrammi complessi; sintassi non intuitive per principianti; dipendenza da Java runtime. \\
\midrule

Draw.io / Lucidchart & GUI visuale; collaborazione in tempo reale; export PNG/SVG. & File binari non versionabili facilmente; diff impossibile; dipendenza da piattaforma cloud. \\
\midrule

Mermaid & Sintasse markdown-friendly; rendering nel browser; integrato in GitHub. & Supporto UML limitato; nessun class diagram completo; meno PlantUML per diagrammi complessi. \\
\midrule

Hand-drawn (disegno manuale) & Massima flessibilit\`{a} visiva. & Non versionabile; non rigenerabile; impossibile da mantenere; non standard. \\
\bottomrule
\end{longtable}

\subsection*{Motivazione}
PlantUML \`{e} lo standard per diagrammi UML in formato testuale: i sorgenti sono versionabili con Git (i diff mostrano le modifiche testuali), rigenerabili in modo deterministic, e supportati dall'estensione VS Code. Draw.io non sarebbe versionabile; Mermaid non supporta i class diagram in modo completo.

\subsection*{Conseguenze}
I 21 sorgenti \texttt{.puml} in \texttt{documentazione/uml/} generano le 25 immagini PNG in \texttt{documentazione/img/}. Ogni modifica ai diagrammi richiede di rigenerare le immagini e ricommetterle insieme ai sorgenti. I sorgenti usano solo costrutti PlantUML standard.

\end{document}
